\documentclass[letterpaper]{article} % DO NOT CHANGE THIS
\usepackage[submission]{aaai2027}  % DO NOT CHANGE THIS
\usepackage[hyphens]{url}  % DO NOT CHANGE THIS
\usepackage{graphicx} % DO NOT CHANGE THIS
\urlstyle{rm} % DO NOT CHANGE THIS
\def\UrlFont{\rm}  % DO NOT CHANGE THIS
\usepackage{natbib}  % DO NOT CHANGE THIS AND DO NOT ADD ANY OPTIONS TO IT
\usepackage{caption} % DO NOT CHANGE THIS AND DO NOT ADD ANY OPTIONS TO IT
\frenchspacing  % DO NOT CHANGE THIS
\usepackage{booktabs}
\pdfinfo{
/TemplateVersion (2027.1)
}

\setcounter{secnumdepth}{2}

\title{LLM-Compiled Wiki vs.\ Vector RAG for Classroom Question Answering}
\author{Anonymous Submission}
\affiliations{Draft --- Sections 1--5}

\begin{document}

\maketitle

\begin{abstract}
Large language models are increasingly deployed as course-specific tutors to enhance student learning, but their usefulness depends on being grounded in vetted instructional materials that are often updated during the semester. Our prior work built a multimodal retrieval-augmented generation (RAG) system over an authentic machine learning course corpus (DS3001 Foundations of Machine Learning) and found that retrieval improved contextual grounding, but that fixed retrieval strategies were suboptimal: performance depended sharply on question specificity and on how text and image context were composed. That result motivates a different question whether how a corpus is structured at ingest time matters more than how much is retrieved at query time. We present a controlled head-to-head comparison of two knowledge representations over an identical classroom corpus: (A) vector RAG, replicating the best-performing configuration from our prior study, and (B) an LLM-compiled wiki, (Karpathy Framework) in which the corpus is synthesized at ingest into linked concept pages with explicit cross-references and citations back to source materials, which the model navigates at question time. We evaluate 59 questions spanning single-fact recall, cross-unit concept linking, synthesis and explanation, and currency after a syllabus revision, scored by LLM judge with a human-authored rubric. Both representations answered single-fact questions about equally well (9.33 vs. 9.96 out of 10), but they diverged sharply on questions that required linking material across course units: the compiled wiki stayed accurate and grounded (9.93, 100% grounded in cited source material) while retrieval's answers both scored lower and were markedly less grounded (8.14, 64% grounded). The compiled representation's citations let a student and instructor trace any claim back to the lecture that introduced it, a critical element for learning. For pedagogical practice, the results indicate a clear advantage for the Wiki based structure such that instructors when using AI methods to enhance student learning should bias toward this structure as it has clear advantages for supporting foundational elements of best practices in learning. 
\end{abstract}

\section{Introduction}

\subsection{Motivation}

Large language models are rapidly being adopted as course-specific tutoring tools. Left ungrounded, however, a general-purpose LLM answers from parametric knowledge that may conflict with how a particular course defines, sequences, or scopes its content; effective educational deployment therefore requires grounding responses in vetted course materials so that answers align with the local curriculum and instructional level \citep{li2025retrieval,jain2025aligning}.

Retrieval-augmented generation (RAG) is the standard mechanism for such grounding \citep{lewis2020retrieval}. In our prior work, we built a multimodal RAG system over authentic materials from a data science course (DS3001)---lecture slides, lecture audio transcripts, and course readings---and found that retrieval substantially improved contextual grounding \citep{wright2026multimodal}. But the same study surfaced a more troubling result: \emph{fixed retrieval strategies are suboptimal}. Performance depended sharply on question specificity and on how the context window was composed. Swapping retrieved text for top-ranked images achieved perfect context recall on specific questions under a five-text, five-image configuration, yet the same strategy degraded faithfulness and factual correctness on generic questions. No single retrieval configuration served all question types well.

This finding reframes the design problem. Vector RAG makes its structural commitments at query time: it retrieves isolated, similarity-ranked chunks and asks the generator to assemble them into a coherent answer. An alternative is to make those commitments at \emph{ingest} time---compiling the corpus into a structured, linked wiki of concept pages, with explicit cross-references and citations back to the raw source materials, which the model then navigates when answering. This ``LLM-compiled wiki'' approach trades query-time chunk retrieval for curated synthesis, explicit concept links, and built-in provenance. If how context is structured and selected matters more than how much is retrieved, then the representation itself---not the retrieval policy---may be the more consequential design choice.

\subsection{Effective Teaching Practices with Chatbots}

Grounding a chatbot in course materials is necessary but not sufficient for it to teach well; a growing body of classroom-deployment research argues that how a tutor is designed to interact matters as much as what it can retrieve. Scaffolded, Socratic designs that pose guiding questions rather than resolving a problem outright have been shown, in a deployed undergraduate CS tutoring system, to shift students from vague help-seeking toward more deliberate, strategic use of the tool \citep{sunil2025socraticai}; teacher presence in how a chatbot-based activity is designed similarly shapes how much students engage with it \citep{li2025teacher}. These effects, however, do not transfer cleanly from controlled studies to real classrooms: an analysis of thousands of real student-tutor conversations found a persistent mismatch between the scaffolding behavior benchmarks assume and how students actually steer real interactions toward their own goals \citep{neagu2026rethinking}---an argument for evaluating on authentic classroom questions over authentic course material, as we do here, rather than on synthetic benchmark sets. Practitioner-derived guidance converges on a related point: generative AI should be taught and used in ways that promote discernment and critical thinking rather than substituting for them \citep{wall2025generative}, and a broader review of classroom deployments finds outcomes depend heavily on how the tool is embedded in instruction rather than on model capability alone \citep{wang2025generative}. Recent field studies of instructors building course-specific AI tutors at institutional scale report that turning these commitments into practice is difficult, and that instructors' biggest challenges center on getting a tutor to reflect a course's own structure and sequencing rather than just answering questions correctly in isolation \citep{ko2026toward}.

That last challenge is where knowledge representation stops being a retrieval-engineering detail and becomes a pedagogical one. A tutor that scaffolds understanding, cites its sources, and helps a student see how today's material builds on an earlier lecture needs access to a course's connective structure---its definitions, its sequencing, its cross-references---not just to the passage nearest a query. This is exactly the structure an LLM-compiled wiki makes explicit at ingest time, while vector RAG reconstructs it, if at all, only implicitly through whatever happens to be retrieved. The comparison in this paper is designed to test whether that difference in representation actually shows up on the question type these teaching practices care about most: not isolated recall, but questions that require connecting ideas across a semester's worth of material.

\subsection{The Gap}

Head-to-head comparisons of compiled wikis against vector RAG are beginning to appear: \citet{cochran2026vector} preregistered such a comparison on a small multi-domain research corpus, finding that vector RAG won on single-fact lookup while the wiki won on cross-paper synthesis, with no architecture dominating every endpoint. But no such comparison exists in an educational setting with authentic classroom data, where the corpus is multimodal, the question distribution is pedagogically shaped, and alignment with vetted materials matters as much as raw correctness \citep{jain2025aligning}. Meanwhile, the graph- and structure-augmented retrieval literature suggests that structure helps multi-hop and relational questions but can hurt simple fact lookup \citep{peng2024graph,zhou2025indepth,aboulela2025exploring}---a trade-off that remains untested on course corpora.

\subsection{Research Questions}

We ask three questions. \textbf{RQ1:} Does an LLM-compiled wiki improve response quality over vector RAG on classroom question answering, holding the generator LLM, corpus, and question set constant? \textbf{RQ2:} How do gains vary by question type---single-fact lookup vs.\ multi-hop concept linking vs.\ synthesis and explanation---extending the generic/specific split of our prior study into a finer taxonomy? \textbf{RQ3:} What are the trade-offs in citation support, robustness to curriculum updates, and ingest/query cost?

\subsection{Contributions}

This paper makes three contributions: (1) the first wiki-vs-RAG comparison on an authentic course corpus, using a text-only subset of DS3001; (2) a fully specified evaluation protocol---a five-part question taxonomy and multi-endpoint metric suite extending our prior 30-question generic/specific set (Section 3)---for validating the comparison at full multimodal scale, of which the present study implements a scoped-down two-class instance; and (3) directional empirical evidence that structure-at-ingest most benefits cross-topic synthesis questions, including a groundedness-without-retrieval-failure finding that complicates the usual retrieval-centric account of RAG's shortcomings.

\section{Background}

\subsection{Learning Outcomes from AI Chatbots}

The case for classroom chatbots ultimately rests on whether they improve learning, not just whether they answer accurately. The experimental evidence here is larger and more consistent than the accuracy literature alone: a meta-analysis of 24 randomized studies found AI chatbots produced a large effect on students' learning outcomes, with stronger effects in higher education than in K--12 settings and, notably, stronger effects for shorter interventions than sustained ones \citep{wu2024chatbots}. A more recent meta-analysis specific to ChatGPT, pooling 35 experimental studies, found a moderate-to-large effect on both cognitive and non-cognitive outcomes, with the instructional mode a chatbot is embedded in---not merely its presence---among the significant moderators \citep{wu2026chatgpt}. Causal evidence points the same direction: a randomized controlled trial in an authentic physics classroom found AI-tutored students outperformed those in active-learning in-class sessions on the same material \citep{kestin2025tutoring}. Retrieval-augmented chatbots specifically, not just chatbots in general, show this pattern too---a RAG-grounded course chatbot deployed in a materials science course was found to measurably enhance students' learning, not merely their perception of it \citep{thway2025harnessing}.

Two things follow from this literature for the present study. First, the moderators of chatbot effectiveness are consistently about \emph{design}---instructional mode, intervention structure, how the tool is embedded---rather than raw model capability, reinforcing the teaching-practices argument of Section 1.2: knowledge representation is one such design choice, and its effect on outcomes is testable rather than assumed. Second, grounding a chatbot in course material is not an end in itself; it is a means to protect the learning-outcome gains this literature documents, since an ungrounded or fabricating chatbot risks reproducing the curriculum misalignment that motivated our prior RAG work in the first place \citep{wright2026multimodal}. The rest of this section reviews the retrieval mechanisms available for that grounding, before turning in Section 2.4 to the education-specific evidence on how well they preserve outcomes when applied to authentic course content.

\subsection{Retrieval-Augmented Generation}

RAG conditions a generator on documents retrieved from an external corpus, originally formulated for knowledge-intensive NLP tasks \citep{lewis2020retrieval}. It is now widely framed as a practical reliability layer: grounding outputs in external evidence, reducing hallucination, and enabling knowledge updates without retraining \citep{peng2024graph,neha2025retrieval,li2024enhancing}. Domain results can be striking---a clinical GPT-4 RAG system reached 96.4\% accuracy with no observed hallucinations \citep{ke2025retrieval}, and the MEGA-RAG pipeline outperformed both LLM-only and standard RAG baselines in public health question answering \citep{xu2025megarag}. Yet the gains are contingent: systematic reviews and unified-framework analyses consistently find that RAG's benefit depends on retrieval relevance, task type, and evaluation design \citep{zhou2025indepth,brown2025systematic}. These contingencies are precisely what our prior study observed at the level of retrieval configuration, and they motivate treating the knowledge representation as an experimental variable in its own right. For a classroom deployment the stakes of that contingency are the outcome gains reviewed in Section 2.1: a RAG system that scores well on its own evaluation metric can still fail to translate into better learning if its retrieval strategy is mismatched to how students actually ask questions.

\subsection{Structured and Graph-Based Retrieval}

Graph-enhanced RAG augments or replaces chunk retrieval with explicit entity and relation structure, supporting multi-hop reasoning and improving explainability \citep{peng2024graph,zhu2025graphsurvey}. The empirical evidence, however, is mixed in an instructive way. KG\textsuperscript{2}RAG improves answer and retrieval F1 on HotpotQA over hybrid RAG baselines \citep{zhu2025kg2rag}; knowledge-graph-extended RAG improves multi-hop MetaQA accuracy at the cost of a slight single-hop degradation \citep{linders2025knowledge}; and graph RAG has been found to \emph{underperform} na\"ive RAG on SQuAD-V2 and TriviaQA even while handling complex relational queries better \citep{aboulela2025exploring}. The recurring pattern---structure helps composition, hurts lookup---is the central trade-off our design tests in the classroom setting.

LLM-compiled wikis are a distinct point in this design space: rather than extracting a graph over the original chunks, the corpus is rewritten at ingest into human-readable, linked concept pages with citations to source material. In the only controlled comparison to date, \citet{cochran2026vector} found vector RAG better at single-fact lookup and the wiki better at cross-paper synthesis, with no single architecture winning every dimension.

\subsection{RAG in Education}

Education-specific work adds a constraint that generic benchmarks miss: retrieval should align to vetted course materials and the local curriculum, not merely to generally correct information \citep{li2025retrieval,jain2025aligning}. The relationship between grounding and quality is also non-monotonic---in math tutoring, humans preferred RAG-grounded responses \emph{unless} grounding became so rigid that helpfulness suffered \citep{levonian2023retrieval}. On the structured side, an educational knowledge-graph plus agentic-RAG system achieved 91.4\% retrieval accuracy with improved learner satisfaction \citep{gao2025karag}, and hybrid RAG with in-context learning outperformed baselines for educational question generation \citep{maity2024leveraging}. Together these results suggest structure has pedagogical value, but no study has isolated the representation itself on an authentic course corpus.

\subsection{Prior Work: Multimodal RAG on DS3001}

This study extends our prior evaluation of multimodal retrieval over DS3001 course data \citep{wright2026multimodal}. That work showed that multimodal retrieval improved context recall, and that selectively replacing retrieved text with top-ranked images helped specific questions---reaching perfect context recall at five text plus five image chunks---while hurting faithfulness and factual correctness on generic questions. The takeaway that motivates the present study is that \emph{how context is structured and selected matters more than how much is retrieved}. If context composition dominates context volume, the natural next experiment compares an ingest-time structured representation against query-time chunk retrieval under matched budgets.

\subsection{Evaluating RAG Systems}

Lexical overlap metrics such as BLEU and ROUGE are unstable under LLM output stochasticity, motivating LLM-as-judge and embedding-based evaluation \citep{lyu2024crudrag}. We retain the RAGAS framework \citep{ragas2023} used in our prior paper for continuity, and follow the multi-endpoint reporting of \citet{cochran2026vector}---scoring accuracy, synthesis quality, citation alignment, and cost separately---to avoid winner-take-all conclusions that single aggregate scores invite \citep{lyu2024crudrag}.

\section{Study Design}

We run a controlled comparison: same corpus, same base LLM, same prompts, same question set; only the knowledge representation and the way context reaches the generator differ. The design is modeled on the preregistered comparison of \citet{cochran2026vector}, combined with the classroom-focused evaluation approach of \citet{jain2025aligning}, including its knowledge-shift testing.

\subsection{Corpus}

The source corpus is the DS3001/DS-3021 course material used in our prior study, together with the instructor's companion online text, the \emph{Foundations of Machine Learning} book (MLBook).\footnote{\url{https://novavolunteer.github.io/MLBook/}, ingested at commit \texttt{40b4199} (2026-05-21).} The study covers the whole course as taught, in 16 topics: course overview and syllabus, development environment, Python basics, ML foundations and the data science lifecycle, data preparation, linear regression, logistic regression, K-Nearest Neighbors, model evaluation metrics, decision trees, decision tree regression, K-means clustering, principal component analysis, ensemble methods/random forest, ML bias and fairness, and the final project. Three legacy course folders that are not part of the course as taught (reinforcement learning, text mining, and a Git primer) are excluded. The \emph{raw corpus} is every text-bearing source file for those topics: 81 documents (56 from the course folder, 25 from the MLBook), about 235k tokens. These are lecture slides (PDF and PowerPoint, including speaker notes), lecture and lab notebooks (markdown and code cells, without outputs), R-to-Python conversion notes, scripts, assignment sheets, the syllabus, and assigned readings. Exact duplicates across the two sources were ingested once. Where both sources held different versions of the same notebook, the current published MLBook version was used, and superseded slide decks and format-only duplicates (PDF exports of included decks, \texttt{.py} exports of included notebooks) were dropped. Every inclusion and exclusion is recorded in the ingest manifest. Text is extracted once, and the same extracted text feeds A-ret and A-full. Data files, spreadsheets, images, and one legacy binary \texttt{.ppt} deck are excluded, consistent with the text-only scope. No lecture transcripts exist for these topics. The \emph{compiled wiki} is the same material synthesized at ingest by an LLM into one concept page per topic (16 pages, about 71k tokens), with cross-links between related pages and citations from every page back to the raw sources, plus navigational and meta pages (\texttt{index.md}, \texttt{getting-started.md}, \texttt{status.md}, and course-level overviews). Both are built from a single frozen snapshot of the source material.

\subsection{Conditions}

We evaluate five conditions (Table~\ref{tab:conditions}). They are designed so that each comparison changes one thing at a time.

\begin{table}[t]
\centering
\small
\setlength{\tabcolsep}{3pt}
\begin{tabular}{@{}l p{1.55cm} p{2.2cm} p{2.0cm}@{}}
\toprule
 & Material & Context selection & Budget \\
\midrule
C0 & none & --- & question only \\
A-ret & raw corpus & top-$k$ vector retrieval & $k{\times}$chunk \\
A-full & raw corpus & everything & unbounded \\
B-nav & wiki & agent reads pages from \texttt{index.md} & matched to A-ret \\
B-full & wiki & everything & unbounded \\
\bottomrule
\end{tabular}
\caption{Evaluation conditions. A-ret vs.\ B-nav compares representations under a matched budget; A-full vs.\ B-full compares them with no selection step; C0 is the closed-book floor.}
\label{tab:conditions}
\end{table}

\textbf{C0 --- Closed book.} The answerer receives the question and no course material, and answers from its parametric knowledge alone. C0 is the floor against which every grounded condition is measured: a grounded condition that does not beat C0 on a question type is not adding course-specific value there, however well it scores in absolute terms.

\textbf{A-ret --- Vector RAG over raw materials.} A replication of our prior pipeline over the raw corpus: text chunked at 1{,}500 tokens with 100-token overlap, embedded with all-mpnet-base-v2,\footnote{all-mpnet-base-v2 has a maximum input length of 384 word-pieces, so each 1{,}500-token chunk is embedded from its first 384 word-pieces only; the generator still receives the full chunk text. Retrieval is therefore driven by the head of each chunk. We keep this configuration for continuity with our prior study rather than re-tuning it, and treat it as a limitation of A-ret rather than of vector RAG in general.} indexed in Pinecone. Token counts use the embedding model's own tokenizer. Chunks never cross document boundaries, and each chunk is passed to the generator as the verbatim source text. The raw corpus yields 203 chunks. For each question, the question is embedded and the $k$ chunks with the highest cosine similarity to it are passed to the generator. $k$ counts chunks, not topics or documents: the $k$ chunks may come from any files and any topics, several may come from the same file, and nothing guarantees that every topic a cross-page question needs is represented. Chunking, embedding model, and index are the best-performing configuration from the prior study, which makes A-ret a strong control rather than a straw man. The prior study used $k{=}5$ over an 80-chunk corpus, about 6\% of the corpus per question. We use $k{=}10$ here, about 5\% of the 203-chunk corpus, to keep the share of the corpus retrievable per question roughly constant as the corpus grew. $k$ also sets B-nav's budget (below), and the choice is explained there. Only content documents are indexed. Navigational and meta files (indexes, status and getting-started pages, READMEs, course overviews) are excluded, because they mention every topic by name, attract retrieval on many queries, and contain no answer content.

\textbf{A-full --- Raw materials in full context.} The entire raw corpus is placed in the answerer's context, with no retrieval step. This is the long-context control. B-full also places all of its material in context, so a B-full advantage over A-ret alone could come from context volume rather than from compilation. A-full vs.\ B-full holds ``everything in context'' fixed and varies only whether that material was compiled.

\textbf{B-nav --- Agentic wiki navigation.} The answerer starts with \texttt{index.md} in context and is given a single tool, \texttt{read\_page(path)}, that returns the text of a wiki page. It may follow the index and the cross-links in each page it reads. Every page read counts against a per-question token budget equal to A-ret's retrieved-context budget ($k \times 1{,}500 = 15{,}000$ tokens), counted with the same tokenizer as A-ret's chunks. A read that would exceed the remaining budget is truncated. Once the budget is exhausted, or after 12 model turns, tool use is disabled and the answerer must answer from what it has read. Requests for pages that do not exist cost nothing, and neither do repeat reads. B-nav can still fail to find the page it needs, because it must choose which pages to read.

\textbf{Choice of budget.} Matching the two budgets keeps A-ret vs.\ B-nav a comparison of representations rather than of context volume, so the budget is raised for both conditions or for neither. The design intent is that B-nav can read the index plus two to three concept pages. In the full-course wiki the index is 765 tokens and concept pages have a median length of about 3{,}900 tokens (range 1{,}300--7{,}500). At the prior study's $k{=}5$ (7{,}500 tokens), B-nav could read only the index plus one or two pages in full, so a cross-page question could exhaust the budget before its second page was read. At $k{=}10$ (15{,}000 tokens), the index plus three median pages fits, and so does the index plus any two concept pages, with the single exception of the two longest pages, of which the last ${\sim}$90 tokens would be cut. We fixed $k{=}10$ before any full-course results were available. We did not shorten or split pages to fit the budget, because page length reflects how much source material a topic has, and splitting pages to suit the evaluation would tune the representation to the test. As a sensitivity analysis, A-ret and B-nav are also run at $k{=}5$ (7{,}500 tokens), to show whether the gap between them depends on the budget. This is the condition the wiki approach actually proposes: the generator navigates the compiled structure instead of being handed all of it. A-ret vs.\ B-nav is therefore the primary test of representation under a matched budget.

\textbf{B-full --- Whole wiki in context.} All wiki pages are placed in the answerer's context. B-full has no selection step, so it cannot fail to find a page, and it serves as the upper bound for the wiki representation. The gap between B-full and B-nav measures how much is lost to navigation.

Held constant across all five conditions: the answerer model, decoding parameters, question set, and judge. The system prompt is shared except for the material-specific instruction. C0 is told to answer from its own knowledge. The other conditions are told to answer only from the provided (or retrieved, or read) course material, cite the source of each claim, and say so if the material does not contain the answer.

\subsection{Questions}

The full protocol extends the prior 30-question generic/specific set into a five-part taxonomy: (a) single-fact lookup, (b) multi-hop concept linking, (c) thematic explanation and synthesis, (d) contradiction and ambiguity handling \citep{hou2024wikicontradict}, and (e) curriculum-update questions probing sensitivity to revised materials. Questions are human-generated with reference answers and word limits, balanced across course topics. The present study uses a coarser instance of this taxonomy with 100 questions:
\begin{itemize}
\item 45 \emph{topic-page} questions answerable from one topic's material, standing in for single-fact lookup;
\item 45 \emph{cross-page} questions requiring material from two or, in a few cases, three topics, standing in for multi-hop and synthesis;
\item 10 \emph{unanswerable} controls: in-scope questions that the course material does not answer.
\end{itemize}
The set was built in two halves. An earlier 100-question set covered only seven of the 16 topics (KNN, decision trees, evaluation metrics, decision tree regression, K-means, ensembles, and bias). A random half of it, stratified by class (23 topic-page, 22 cross-page, 5 unanswerable), was retained. The other 50 questions were newly generated for the nine added topics with the same procedure: 22 topic-page questions (two or three per new topic); 23 cross-page questions, each over a randomly drawn pair (21) or triple (2) of topics that includes at least one new topic, so that links between new and original material are tested; and 5 unanswerable controls. Retained unanswerable controls were re-checked against the expanded corpus.
The questions were generated from the \emph{raw corpus only}; the generator never saw the compiled wiki. This keeps the wiki's structure and wording from shaping the questions. Generation used a model from a different provider than the answerer (\texttt{gpt-5.5}), so that the answerer is not answering questions written by its own model family. Each topic-page question was generated from that topic's files. Each cross-page question was generated from the combined files of its topics, and was kept only if its cited sources spanned all of them. Every question comes with a reference answer of at most 150 words, written from the material, and with its \emph{gold source}: the raw files the reference answer relies on. Each gold file maps to its topic's concept page, which is the gold unit for the wiki conditions. Each unanswerable candidate was checked against the full raw corpus and discarded if any file answered it. Some of the retained controls name a topic that the material mentions without explaining, such as the gap statistic. A question can draw on detail that the wiki omitted; the wiki conditions then lose that question, and that loss is a real cost of compiling the material. The course instructor reviewed every generated question and reference answer, and edited them where needed.

\subsection{Logging}

For every condition and question we log three quantities, so that differences in answer quality can be attributed to what the answerer actually saw.
\begin{itemize}
\item \textbf{Input tokens}, reported per condition as mean, standard deviation, and range across questions. This is the answerer's total input on its final call: instructions, material, and question. For B-nav this includes the index and every page read. We separately log the tokens billed across all of B-nav's calls, and the course-material tokens alone, counted with the budget tokenizer.
\item \textbf{Units seen}: the chunk IDs retrieved (A-ret) or pages read, in order (B-nav). For the full-context conditions this is the whole corpus or wiki.
\item \textbf{Gold source in context}: whether the material needed to answer was in the answerer's context. For topic-page questions this is binary. For cross-page questions we log the fraction of required sources present. For B-nav and B-full the unit is a concept page. For A-ret and A-full it is any chunk from a raw document cited by the gold page. This is trivially true for the full-context conditions and false for C0, and is the variable that separates retrieval or navigation failures from generation failures.
\end{itemize}

\subsection{Scoring}

An LLM judge from a different provider than the answerer (\texttt{gpt-5-mini}), so that the model is not grading its own output, grades each response, with condition labels hidden. The judge's reference is built the same way for every condition: the full compiled wiki; the question's gold raw files; every raw file the response cites; for A-ret, the retrieved excerpts; and the question's reference answer. A claim taken correctly from a slide or notebook is therefore credited even when the wiki does not repeat it, provided the response cites that file or it is a gold source. The full wiki plus the full raw corpus (about 306k tokens) no longer fits reliably in the judge's context window, which is why the judge does not see the whole raw corpus. In place of the RAGAS triad \citep{ragas2023} retained in the full protocol, the judge returns a 1--10 correctness score and a binary groundedness flag: whether every substantive claim in the answer is stated in the course material. These are lightweight proxies for RAGAS factual correctness and faithfulness respectively, not substitutes for them. Groundedness is judged against the course material, not against general ML knowledge. For C0, an answer that is correct but drawn from general knowledge is therefore marked ungrounded, and C0's fabrication rate is an upper bound on its true rate of fabrication.

\textbf{Refusals.} Under the previous rubric, a correct ``this is not in my materials'' scored the same as a fabricated answer. We fix the following rule before running any of the new conditions. The judge first classifies each response as an \emph{answer} or a \emph{refusal}, meaning the response states that the provided material does not contain the answer, or otherwise declines to answer. Refusals receive no 1--10 score. Each refusal is labelled \emph{justified} if any of the question's gold sources was missing from the answerer's context, and \emph{unjustified} if all of them were present, using the logged source-in-context coverage. A response that answers part of the question and says the rest is not in the material counts as an answer and is scored. The judge also records whether an answer explicitly tells the student that some of what it says is not covered by the course material. An ungrounded answer that says so is \emph{flagged}. Only an ungrounded answer that does not say so counts as \emph{fabrication}, because a student cannot tell that it goes beyond the course. Every question therefore ends in exactly one of five outcomes: grounded answer, flagged ungrounded answer, fabrication, justified refusal, or unjustified refusal. On an unanswerable control, a correct response says that the material does not cover the question. If it then offers only related course content, it is a grounded answer. We report, per condition and question class:
\begin{itemize}
\item the mean correctness score over answered questions;
\item the answer rate;
\item the fabrication and flagged-ungrounded rates over all questions;
\item the justified and unjustified refusal rates.
\end{itemize}
As a single summary, we also report mean correctness over all questions with every refusal scored at the floor (1). This summary penalizes refusals and never rewards them, so it cannot favor a condition that refuses often. A justified refusal is never counted as fabrication, and an unjustified refusal is counted as a miss. C0 has no material to refer to, so its refusals are reported but not split into justified and unjustified.

\subsection{Statistical Analysis and Full Protocol}

We report 95\% bootstrap confidence intervals by resampling questions. Pairwise differences between conditions are computed per question and bootstrapped as paired differences. The primary analysis uses the $k{=}10$ / 15{,}000-token budget. The $k{=}5$ / 7{,}500-token run of A-ret and B-nav is reported as a secondary analysis of the same matched-budget contrast. The full protocol additionally repeats each run (10 rounds, as in our prior paper) to separate answerer and judge stochasticity from question-sampling error. It also adds claim--citation alignment, cost (ingest compute, query tokens, latency), and a knowledge-shift phase in which a subset of source materials is revised, both representations are re-ingested, and affected questions are re-asked to measure staleness and update cost \citep{jain2025aligning,hou2024wikicontradict}.

\subsection{Hypotheses}

We preregister three hypotheses, stated over the matched-budget pair. \textbf{H1:} A-ret performs at least as well as B-nav on single-fact (topic-page) questions, at lower cost. \textbf{H2:} B-nav outperforms A-ret on multi-hop and synthesis (cross-page) questions. \textbf{H3:} across endpoints the outcome is a set of trade-offs rather than a single winner, consistent with \citet{cochran2026vector}. We treat the long-context pair (A-full vs.\ B-full) and the navigation gap (B-full vs.\ B-nav) as secondary analyses, and every condition is compared against C0. A fourth candidate hypothesis, that the wiki gives better claim--citation alignment and provenance, is not preregistered because no endpoint in this study measures citation alignment. The closest available proxy, the binary groundedness flag, is too coarse to stand in for it.

\subsection{What Has Been Run}

A pilot run, reported in Section~4, executed B-full and a vector-retrieval condition over the \emph{compiled wiki pages} rather than over the raw corpus. That condition indexed all 12 wiki pages, navigational pages included, into 21 vectors. It used the earlier rubric, under which refusals were not distinguished, and it did not apply all of the logging described above. Its chunker also rebuilt each chunk from embedding-tokenizer IDs, so its answerer received lowercased text with altered spacing (e.g.\ ``gini = 0. 5''). This alone may have depressed the pilot's retrieval scores. The pilot's retrieval condition is superseded by A-ret, which passes verbatim text. All five conditions are re-run under the protocol above, with B-full included, so that every condition shares the same rubric, judge reference, and logging. A first run of this protocol over the seven-topic corpus was superseded before completion when the corpus was expanded to the full course. Its partial results are archived and not reported.


\section{Results}

Given time and resource constraints, we were not able to execute the full multimodal, RAGAS-scored protocol described in Section 3 on the complete DS3001 corpus. The two-arm evaluation reported in this section is this paper's empirical contribution: it validates the pipeline end-to-end and gives a directional read on H1--H2. It is scoped down from the full design in ways we flag explicitly in Section 4.4: text-only rather than multimodal, scored with a single LLM-judge rating rather than the RAGAS triad, and each arm run only once (bootstrapped over questions for the CIs below, but not over repeated answerer/judge sampling). We report it as informative on its own terms, not as a substitute for RQ1--RQ3 answered at full scale, which we leave to future work (Section 5.4).

\subsection{Headline Results}

\begin{table}[t]
\centering
\small
\begin{tabular}{lcc}
\toprule
 & Arm A (RAG) & Arm B (Wiki) \\
\midrule
Avg.\ score (1--10) & 9.05 [8.49, 9.54] & \textbf{9.95} [9.88, 10.00] \\
Grounded rate & 81\% [71\%, 90\%] & \textbf{98\%} [95\%, 100\%] \\
\bottomrule
\end{tabular}
\caption{Headline results with 95\% bootstrap CIs (question resampling, $B{=}20{,}000$), 59 questions, both arms graded 59/59 with zero errors.}
\label{tab:headline}
\end{table}

Table~\ref{tab:headline} shows the wiki ahead of vector RAG on both measures at this scale, with the larger gap in groundedness (17 points, 95\% CI on the difference [7, 27] points) than in raw score (0.90 points, 95\% CI [0.42, 1.44])---both CIs exclude zero---consistent with the qualitative pattern that RAG can produce a fluent, plausible-sounding answer that leans on knowledge outside what it actually retrieved.

\begin{table}[t]
\centering
\small
\setlength{\tabcolsep}{4pt}
\begin{tabular}{lcccc}
\toprule
 & \multicolumn{2}{c}{Arm A (RAG)} & \multicolumn{2}{c}{Arm B (Wiki)} \\
 & Score & Grounded & Score & Grounded \\
\midrule
Topic-page ($n{=}45$) & 9.33 & 87\% & 9.96 & 98\% \\
Cross-page ($n{=}14$) & 8.14 & 64\% & 9.93 & 100\% \\
\bottomrule
\end{tabular}
\caption{Results split by question class. The wiki's advantage over RAG is roughly three times larger on cross-page questions than on topic-page questions, in both score and groundedness.}
\label{tab:byclass}
\end{table}

Table~\ref{tab:byclass} is the more informative result: splitting by question class reproduces, directionally, the pattern \citet{cochran2026vector} report on their research-paper corpus---RAG is closer to competitive on single-topic questions and falls further behind on synthesis. Full-context answering (Arm B) is nearly indifferent to question class (9.96 vs.\ 9.93 score, 98\% vs.\ 100\% grounded); vector RAG is not (9.33 vs.\ 8.14 score, an 0.62-point gap on topic-page [95\% CI on the gap: 0.16, 1.18] widening to 1.79 on cross-page [95\% CI: 0.57, 3.14]; 87\% vs.\ 64\% grounded). The score gap is statistically distinguishable from zero in both classes, but the two classes' groundedness gaps are not equally well established: the cross-page grounded-rate gap (36 points, 95\% CI [14, 64]) clearly excludes zero at $n=14$, while the topic-page grounded-rate gap (11 points, 95\% CI [0, 22]) just touches zero at $n=45$---the one comparison in this study that does not reach significance at the 95\% level.

\subsection{Where Vector RAG Lost Ground}

Retrieval itself was largely successful: 43 of 45 topic-page questions (96\%) retrieved a chunk from the question's actual source page in the top-5. The two misses both pulled in short navigational pages (index, getting-started, status, course overview) instead of the target concept page, on questions about ML Bias interventions and about which k-means slide deck is canonical---both scored low (3 and 4) and were marked ungrounded, as expected when the answerer never saw the relevant material.

More striking is that retrieval failure explains only 2 of the 11 ungrounded Arm A answers. The other 9 retrieved the correct page (or, for cross-page questions, relevant chunks from more than one page) and were \emph{still} marked ungrounded---the answerer added unsupported detail even with the right excerpts in context. This suggests chunked context invites fabrication in a way full-document context does not, independent of whether retrieval itself succeeded, and is worth isolating explicitly once the full RAGAS faithfulness metric is in place.

\subsection{Relation to the Preregistered Hypotheses}

These results speak only partially to H1--H3. \textbf{H1} (RAG competitive on single-fact lookup, at lower cost) is not well supported once the bootstrap CIs are considered: Arm A's topic-page score (9.33) is numerically close to Arm B's (9.96), but the 95\% CI on that gap excludes zero (Section~4.1), so the two arms are statistically distinguishable even on single-fact lookup, and we have not yet measured the cost side of H1 at all. The one part of H1 the data does not rule out is groundedness specifically: the topic-page grounded-rate gap (87\% vs.\ 98\%) is the only comparison in this study whose CI touches zero, so RAG's shortfall there is directionally consistent but not statistically established at $n=45$. \textbf{H2} (wiki outperforms on multi-hop/synthesis) is supported, and more strongly than H1 is disconfirmed: the wiki's advantage on cross-page questions is roughly three times its advantage on topic-page questions, and unlike the topic-page groundedness gap, every cross-page comparison's CI excludes zero. \textbf{H3} (trade-offs rather than one architecture dominating) is not supported by these results specifically---the wiki led on both question classes at this corpus scale, and the score gap is statistically distinguishable from zero in both---though whether that holds at full DS3001 scale, with cost and citation endpoints included, is a question for future work (Section 5.4).

\subsection{Scope of This Evaluation}

This evaluation differs from the design in Section 3 in ways that bound how much weight its results should carry: (1) the corpus is text-only and roughly an order of magnitude smaller than the full multimodal DS3001 corpus targeted by the study design; (2) scoring uses a single LLM-judge rating and a binary groundedness flag rather than the full RAGAS context-recall/faithfulness/factual-correctness triad; (3) each arm was answered and graded only once per question rather than repeated; we bootstrap the 59 questions with replacement ($B{=}20{,}000$) to report 95\% CIs on question-sampling error, but this does not capture answerer or judge stochasticity (the same question re-answered, or the same answer re-graded, may score differently), which needs repeated runs rather than resampling of a single run; (4) judging was not blinded to arm; (5) cost (ingest compute, query tokens, latency) was not measured; and (6) the knowledge-shift/curriculum-revision phase has not been run. We report these results as evidence the pipeline works end-to-end and as a directional signal worth following up, not as a substitute for the full protocol.

\section{Discussion}

\subsection{RQ1/RQ2: Structure-at-Ingest vs.\ Retrieval-at-Query}

At the scale of this evaluation, the LLM-compiled wiki answers questions better than vector RAG on both question classes, and the advantage is not uniform: it is roughly three times larger on cross-page synthesis questions than on single-topic lookup, in both correctness score and groundedness (Section~4.1). Bootstrap CIs over the 59 questions (Section~4.1) show this pattern is not just noise from a small question set---every gap except one (the topic-page groundedness gap) is statistically distinguishable from zero---though those CIs speak only to question-sampling error, not to answerer or judge stochasticity, which a single run cannot separate out (Section~4.4). RQ1 therefore has a directional answer---an LLM-compiled wiki improves response quality over vector RAG under a matched generator, prompts, context budget, and question set---and RQ2's answer is that the improvement concentrates in exactly the question type the literature would predict: questions that require linking material across the corpus rather than retrieving a single relevant passage \citep{cochran2026vector,peng2024graph,aboulela2025exploring}. We still stress \emph{directional} over definitive: this is a single run on a 21-vector index, and Section~4.4 lists the further ways it falls short of the protocol in Section~3.

\subsection{Structure Also Reduces Fabrication, Not Just Retrieval Misses}

The most surprising result is in Section~4.2: only 2 of vector RAG's 11 ungrounded answers were retrieval misses. The other 9 retrieved the right material and were still marked ungrounded---the answerer added detail the retrieved excerpts did not support. The graph- and structure-augmented retrieval literature typically frames structure as a fix for \emph{retrieval} (finding the right chunk) rather than \emph{generation} (using only what was found) \citep{peng2024graph,zhu2025graphsurvey}. Our results suggest structure-at-ingest does both: the wiki's cross-referenced concept pages appear to constrain the generator's elaboration as well as its access to relevant content. This reframes the RAG-vs-wiki question---even a hypothetical vector RAG system with perfect retrieval might not close the groundedness gap, because part of the gap looks like an artifact of chunked, decontextualized excerpts inviting a generator to fill in surrounding context on its own, not purely a retrieval-quality problem.

\subsection{Limitations and Threats to Validity}

Beyond the scoping limits noted in Section~4.4, three broader threats apply: (1) generalization is untested beyond a single course corpus (DS3001) and its topic mix; (2) our LLM-judge scores and binary groundedness flag have not been validated against human raters, so judge reliability is an open question; and (3) the chunking/embedding configuration for Arm A is inherited unchanged from our prior study (Section~2.5) to isolate the representation variable, but this means a different RAG configuration is a possible confound---the reported gap could narrow if Arm A were re-tuned for this corpus rather than reused from the prior one.

\subsection{Future Work}

This study leaves RQ3 (citation support, robustness to curriculum updates, ingest/query cost) essentially untested, and several elements of the protocol in Section~3 were not executed: repeated-run bootstrapping to isolate answerer and judge stochasticity from question-sampling error (we bootstrap the latter alone by resampling the fixed set of 59 questions, per Section~4.1), blinded judging, the RAGAS context-recall/faithfulness/factual-correctness triad in place of our single-score proxy, the knowledge-shift phase, and---most consequentially---the complete multimodal DS3001 corpus (lecture slides, audio transcripts, and readings) in place of the 12-page text-only subset used here. A full-scale run following Section~3 as specified would confirm whether the directional pattern observed here (wiki ahead on both classes, more so on synthesis) survives a realistically sized index, and would answer RQ3 by measuring whether the wiki's quality advantage justifies its ingest-time synthesis cost against vector RAG's lower per-query cost. We leave that run, and with it a settled answer to RQ1--RQ3, to future work.

\bibliography{references}

\end{document}
